% Ad Familiares 7.9 --- headnote
% ad-familiares-07-09, mid-September 54 BC, Rome

\textit{Cicero to C. Trebatius Testa, written from Rome in
mid-September 54 BC. Trebatius is still with Caesar's army in
Gaul, where Cicero had placed him on the staff at the start of
the year (the recommendation is \textit{Fam.} 7.5). The
\textit{luctum} that keeps Cicero from writing to Caesar directly
is the death of Caesar's daughter Julia --- Pompey's wife and the
last familial hinge between the two men --- who died in
childbirth in late August or early September. Cicero routes his
solicitude for Trebatius through Balbus instead.}

\textit{The letter is among the shortest in the corpus and runs
on the dry register of the whole Trebatius sequence. The middle
section --- \textit{tu tibi desse noli; serius potius ad nos, dum
plenior} --- presses the running joke that Trebatius should not
flee Gaul empty-handed, however little he is enjoying the camp;
``especially now that Battara is dead'' is private shorthand,
some Roman rival or distraction no longer giving Trebatius any
reason to hurry home. The closing vignette is pure social
comedy: a man whose name Cicero affects not to be able to
remember --- ``a certain Cn. Octavius, or is it Cn. Cornelius?'' ---
and whom he characterises in a single self-cancelling phrase
(\textit{summo genere natus, terrae filius}, ``born of the highest
family, a son of the earth'' --- the second tag is the standard
Latin idiom for a nobody) keeps inviting him to dinner on the
strength of knowing Trebatius. Cicero never goes, but enjoys
the compliment.}
